% Originally: content/week-08.tex had a \section{Solutions} holding nothing but a placeholder
% note, which had already gone stale, since the statements were in fact transcribed. So 8.5
% stood in the document unsolved. Written here for the first time.
\section{Solutions}
\label{sec:change_of_variables_solutions}

\begin{exercisesolution}[Solution to \cref{ex:8.5}]
% Generator: Claude Opus 5 (High)
In each part the work is the same two steps: find $\lvert\det\Jac\rvert$, then translate every
constraint. As the exercise says, no integral is actually evaluated.
\begin{enumerate}[label=\textbf{\arabic*.}]
  \item \textbf{Polar coordinates.} Here $\Phi(r,\theta) := (r\cos\theta, r\sin\theta)$ has
    $\det\Jac\Phi = r > 0$ (\cref{ex:polar_spherical_coordinates}), so
    $dx_1dx_2 = r\,dr\,d\theta$. The three constraints translate one at a time:
    \begin{itemize}
      \item $1 < x_1^2+x_2^2 < 4$ becomes $1 < r < 2$;
      \item $x_1 > 0$ becomes $\cos\theta > 0$, i.e.\ $\theta \in (-\tfrac\pi2, \tfrac\pi2)$;
      \item $x_2 < x_1$ becomes $r\sin\theta < r\cos\theta$; since $r > 0$ and $\cos\theta > 0$
        this is $\tan\theta < 1$, i.e.\ $\theta < \tfrac\pi4$.
    \end{itemize}
    So the transformed domain is the rectangle
    $(r,\theta) \in (1,2)\times\big(-\tfrac\pi2,\tfrac\pi4\big)$ and
    \[\int_A x_1^2\sin(x_2)\,dx_1dx_2
      = \int_{1}^{2}\!\!\int_{-\pi/2}^{\pi/4} r^3\cos^2\theta\,\sin(r\sin\theta)\,d\theta\,dr,\]
    the extra power of $r$ coming from $x_1^2 = r^2\cos^2\theta$ times the Jacobian factor $r$.
  \item \textbf{$u := xy$, $v := x^2$.} Follow Sascha Brack's hint and fix the signs first. On $B$
    the constraint $x^2 < y$ forces $y > 0$, and then $xy > 1 > 0$ forces $x > 0$; so
    $u > 0$ and $v > 0$ throughout, and $x = \sqrt v$, $y = u/\sqrt v$.

    Differentiating in the convenient direction,
    \[\frac{\partial(u,v)}{\partial(x,y)}
      = \det\begin{pmatrix} y & x \\ 2x & 0\end{pmatrix} = -2x^2 = -2v,
      \qquad\text{so}\qquad dx\,dy = \frac{du\,dv}{2v}.\]
    Now the constraints. $1 < xy < 2$ becomes $1 < u < 2$ at once. The second is the interesting
    one: $x^2 < y < 2x^2$ reads $v < u/\sqrt v < 2v$, i.e.
    \[v^{3/2} < u < 2\,v^{3/2}, \qquad\text{equivalently}\qquad
      \Big(\frac{u}{2}\Big)^{2/3} < v < u^{2/3}.\]
    The integrand becomes $y^2e^{-xy} = (u^2/v)\,e^{-u}$, so
    \[\int_B y^2e^{-xy}\,dx\,dy
      = \int_{1}^{2}\!\!\int_{(u/2)^{2/3}}^{u^{2/3}} \frac{u^2e^{-u}}{2v^2}\,dv\,du.\]
    Note what did \emph{not} happen: straightening $1 < xy < 2$ into a pair of constant bounds
    did nothing for the other constraint, which stays curved. Straightening one constraint does
    not oblige the rest to cooperate.
  \item \textbf{A linear change of variables.} With $u := z$, $v := z-2y$, $w := 3x+y+z$,
    \[\frac{\partial(u,v,w)}{\partial(x,y,z)}
      = \det\begin{pmatrix} 0 & 0 & 1 \\ 0 & -2 & 1 \\ 3 & 1 & 1\end{pmatrix} = 6,
      \qquad\text{so}\qquad dx\,dy\,dz = \frac{du\,dv\,dw}{6},\]
    a constant, since the map is linear. Each constraint already \emph{is} one of the new
    coordinates, so the transformed domain is literally the box
    \[(u,v,w) \in (0,1)\times(1,2)\times(-2,0).\]
    Inverting the substitution gives $z = u$, $y = \tfrac{u-v}{2}$ and
    $x = \tfrac{1}{6}\big(2w - 3u + v\big)$, so $xyz = \tfrac{1}{12}\,u\,(u-v)\,(2w-3u+v)$ and
    \[\int_C xyz\,dx\,dy\,dz
      = \int_{0}^{1}\!\!\int_{1}^{2}\!\!\int_{-2}^{0}
        \frac{u\,(u-v)\,(2w-3u+v)}{72}\,dw\,dv\,du.\]
\end{enumerate}

\begin{ainote}
The three parts are graded deliberately, and comparing them is the lesson. In part 3 the change
of variables was \emph{chosen} to be the constraints, so the domain is a box and the Jacobian a
constant, and nothing is left to do. In part 1 the domain is a box in the new coordinates but the
Jacobian is not constant. In part 2 neither holds: the Jacobian varies and one of the two
constraints stays curved. Reading the substitution off the constraints, as in part 3, is the move
worth remembering, and it is available far more often than it looks.
\end{ainote}
\end{exercisesolution}

\begin{exercisesolution}[Solution to \cref{ex:ai_triangle_substitution}]
% Generator: Claude Opus 5 (High)
\textbf{1. The new domain and the Jacobian factor.} The map $\Psi(x,y) := (y-x,\ y+x)$ is
linear and invertible, with
\[\Psi^{-1}(u,v) = (x,y) = \left(\frac{v-u}{2},\ \frac{u+v}{2}\right),
  \qquad
  \det\Jac\Psi^{-1} = \det\begin{pmatrix} -1/2 & 1/2 \\ 1/2 & 1/2\end{pmatrix} = -\frac12,\]
so $dx\,dy = \tfrac12\,du\,dv$. The three sides of $T$ translate one at a time: $x \geq 0$ reads
$u \leq v$, the condition $y \geq 0$ reads $u \geq -v$, and $x+y \leq 1$ reads $v \leq 1$. Taking
strict inequalities for the interior gives exactly $0 < v < 1$ and $-v < u < v$, a triangle
again, but one whose horizontal slices run in the direction the integrand likes.

\textbf{2. The value.} On that region the integrand is $e^{u/v}$, and $v$ is constant along the
inner integration, so the inner integral is an exponential in $u$:
\[I = \frac12\int_0^1\!\!\int_{-v}^{v} e^{u/v}\,du\,dv
    = \frac12\int_0^1 \Big[v\,e^{u/v}\Big]_{u=-v}^{u=v}\,dv
    = \frac{e - e^{-1}}{2}\int_0^1 v\,dv
    = \frac{e - e^{-1}}{4}.\]

\textbf{3. The origin.} A single point is a Jordan null set, so removing it changes neither the
Jordan measurability of $T$ nor the value of the integral. What does need checking is that the
integral exists at all, and it does, because the integrand is \emph{bounded}: on $T$ one has
$\lvert y-x\rvert \leq y+x$, so the exponent $u/v$ lies in $[-1,1]$ and the integrand lies
between $e^{-1}$ and $e$. It is the boundedness, and not the smallness of the bad set on its
own, that rescues the computation.

\begin{remark}[Two ways to choose a substitution]
\label{rem:substitution_from_the_integrand}
\Cref{ex:change_of_variables_domain} chose $\Phi$ to straighten the \emph{domain}. Here the
choice comes from the \emph{integrand}, which is a function of $y-x$ and $y+x$ and of nothing
else. That the domain also becomes pleasant is a bonus, and a mild one: the new region is still
a triangle. What made the integral collapse is that the inner integration now runs along the
direction in which the integrand is a plain exponential.
\end{remark}
\end{exercisesolution}

\begin{exercisesolution}[Solution to \cref{ex:ai_order_swap_gaussian}]
% Generator: Claude Opus 5 (High)
\textbf{1.} The inner integral asks for an antiderivative of $y \mapsto e^{y^2}$, and there is
none in elementary terms. Nothing about the outer integral repairs this, so the order as written
is a dead end.

\textbf{2.} The region is the triangle below the diagonal, and it admits two descriptions:
\[R = \{(x,y) : 0 \leq x \leq 1,\ x \leq y \leq 1\}
    = \{(x,y) : 0 \leq y \leq 1,\ 0 \leq x \leq y\}.\]
The first reads it by vertical slices, the second by horizontal ones. Taking the second and
applying \cref{thm:fubini},
\[I = \int_0^1\!\!\int_0^{y} e^{y^2}\,dx\,dy.\]

\textbf{3.} Now $e^{y^2}$ is constant in $x$, so the inner integral contributes a factor of $y$,
and that factor is exactly what the substitution $s := y^2$ needs:
\[I = \int_0^1 y\,e^{y^2}\,dy = \left[\frac{e^{y^2}}{2}\right]_0^1 = \frac{e-1}{2}.\]

\begin{remark}[What the swap actually bought]
\label{rem:order_swap_supplies_the_factor}
The reversal did not find a cleverer antiderivative. It produced a factor of $y$ that was not
there before, and $y\,e^{y^2}$ integrates where $e^{y^2}$ does not. That is the pattern worth
recognising: when the length of a slice depends on the variable that is causing the trouble,
swapping the order hands you that length as a factor in the integrand.
\end{remark}
\end{exercisesolution}

\begin{exercisesolution}[Solution to \cref{ex:ai_frullani}]
% Generator: Claude Opus 5 (High)
\textbf{1.} Both exponentials equal $1$ at $x = 0$, so the quotient is a difference quotient in
disguise. Expanding, $e^{-ax} - e^{-bx} = (b-a)x + O(x^2)$, and therefore
\[\lim_{x\to0^+}\frac{e^{-ax} - e^{-bx}}{x} = b - a.\]
Assigning that value at $x = 0$ makes the integrand continuous on $[0,\infty)$, and it decays
exponentially, so the improper integral converges.

\textbf{2.} Only the second exponential depends on $t$, and the $x$ that the chain rule produces
cancels the one in the denominator:
\[I'(t) = \int_0^\infty \frac{\partial}{\partial t}\left(\frac{e^{-ax} - e^{-tx}}{x}\right)dx
        = \int_0^\infty \frac{x\,e^{-tx}}{x}\,dx
        = \int_0^\infty e^{-tx}\,dx
        = \frac{1}{t}.\]
That cancellation is the whole trick. Differentiating in the parameter removes the $1/x$ which
made the integrand hard, and leaves an integral of one line.

\textbf{3.} At $t = a$ the two exponentials coincide and the integrand vanishes identically, so
$I(a) = 0$. Consequently, by the fundamental theorem of calculus,
\[\int_0^\infty \frac{e^{-ax} - e^{-bx}}{x}\,dx = I(b) = I(b) - I(a) = \int_a^b \frac{dt}{t}
  = \log\!\left(\frac{b}{a}\right).\]
That the answer sees only the ratio $b/a$ is now visible in the derivation rather than a
surprise, since $I'(t) = 1/t$ is scale-invariant.

\begin{remark}[The hypotheses of \cref{thm:feynmans_trick} do not literally cover this]
\label{rem:feynman_on_a_half_line}
As stated, \cref{thm:feynmans_trick} differentiates under an integral over a \emph{compact} set
$K$, and here the domain is $[0,\infty)$. The step is still legitimate, by the same reasoning
\cref{ex:feynman_dirichlet_integral} appeals to. Fix $0 < t_0 < t_1$. For $t \in [t_0,t_1]$ the
$t$-derivative of the integrand is $e^{-tx} \leq e^{-t_0x}$, a bound that is independent of $t$
and integrable on $[0,\infty)$, so the differentiation may be carried out on $[0,R]$ and the
limit $R \to \infty$ taken afterwards. The compact statement is the one the course proves; the
dominated version is what the applications actually use.
\end{remark}
\end{exercisesolution}

\begin{exercisesolution}[Proof of \cref{ex:volume_ellipsoid_change_variables}]
% Generator: Gemini 3.6 Flash (High)
% Correction: Claude Opus 5 (High) --- \operatorname{vol} -> \vol, J\Phi -> \Jac\Phi, ^T ->
% \transp, and \iiint -> \int as elsewhere in this chapter.
\begin{enumerate}[label=\textbf{(\alph*)}]
  \item Define $\Phi(u,v,w) := (au, bv, cw)\transp$, and write $(x,y,z) = \Phi(u,v,w)$. Then
    $(x/a)^2 + (y/b)^2 + (z/c)^2 = u^2+v^2+w^2$, so $(u,v,w)$ lies in the closed unit ball if and
    only if $(x,y,z) \in U$. Since $a,b,c > 0$, the map $\Phi$ is a linear bijection, and therefore
    $\Phi(\overline{B_1(0)}) = U$.

  \item The Jacobian matrix of $\Phi$ is the diagonal matrix
\[\Jac\Phi = \begin{pmatrix} a & 0 & 0 \\ 0 & b & 0 \\ 0 & 0 & c \end{pmatrix},\]
    constant in $(u,v,w)$, with determinant $\det \Jac\Phi = abc > 0$.

  \item By the change of variables theorem (\cref{thm:change_of_variables}),
\[\vol_3(U) = \int_U 1 \, d\vol_3 = \int_{\overline{B_1(0)}} \lvert\det \Jac\Phi\rvert \, d\vol_3 = abc \int_{\overline{B_1(0)}} 1 \, d\vol_3 .\]
The unit ball in $\mathbb{R}^3$ has volume $\tfrac{4}{3}\pi$, and therefore
\[\vol_3(U) = \frac{4}{3} \pi abc.\]
\end{enumerate}
\end{exercisesolution}

\begin{exercisesolution}[Proof of \cref{ex:cubic_shear_area}]
% Generator: Claude Opus 5 (High)
\begin{enumerate}[label=\textbf{(\alph*)}]
  \item Differentiating each component,
\[D_x\varphi = \Jac\varphi(x) = \begin{pmatrix} 3x_1^2 & -1 \\ 1 & 1 \end{pmatrix},
  \qquad \det \Jac\varphi(x) = 3x_1^2 + 1 \ge 1 > 0 .\]

  \item The determinant never vanishes, so by the inverse function theorem
    (\cref{thm:inverse_function_theorem}) $\varphi$ is a local diffeomorphism at every point; in
    particular $\varphi(\mathbb{R}^2)$ is open. It remains to prove global injectivity, and then
    $\varphi$ is a bijection onto its image which is locally a diffeomorphism, hence a
    diffeomorphism onto its image.

    Suppose $\varphi(x) = \varphi(y)$. The second components give $x_2 + x_1 = y_2 + y_1$, that is,
    $x_2 - y_2 = -(x_1 - y_1)$. Substituting this into the equality of first components,
    $x_1^3 - x_2 = y_1^3 - y_2$, yields
\[x_1^3 - y_1^3 = x_2 - y_2 = -(x_1 - y_1),
  \qquad\text{that is,}\qquad (x_1 - y_1)\big(x_1^2 + x_1y_1 + y_1^2 + 1\big) = 0 .\]
    The second factor is at least $1$, because $x_1^2 + x_1y_1 + y_1^2 = \big(x_1 +
    \tfrac{y_1}{2}\big)^2 + \tfrac34 y_1^2 \ge 0$. Hence $x_1 = y_1$, and then $x_2 = y_2$.

  \item By part \textbf{(a)} the Jacobian determinant is positive, so the change of variables
    theorem (\cref{thm:change_of_variables}) gives
\[\vol_2(\varphi(Q)) = \int_Q \lvert\det \Jac\varphi\rvert \, dx
  = \int_0^1\!\!\int_0^1 \big(3x_1^2 + 1\big) \, dx_2\,dx_1
  = \int_0^1 \big(3x_1^2+1\big)\,dx_1 = 1 + 1 = 2 .\]
\end{enumerate}
\begin{remark}
The map is not linear, so the image of the square is not a parallelogram and there is no single
factor by which areas scale. The change of variables formula handles that by scaling
\emph{pointwise} and integrating the result, which is the whole reason the determinant appears
inside the integral rather than outside it.
\end{remark}
\end{exercisesolution}

\begin{exercisesolution}[Solution of \cref{ex:gaussian_second_moment_polar}]
% Generator: Claude Opus 5 (High)
Following the hint, the substitution $(x,y) \mapsto (y,x)$ is a linear change of variables with
Jacobian determinant $-1$, and it carries the integrand $x^2e^{-\alpha(x^2+y^2)}$ to
$y^2e^{-\alpha(x^2+y^2)}$. Hence the two integrals agree, and adding them gives
\[2I_\alpha = \int_{\mathbb{R}^2} (x^2+y^2)\,e^{-\alpha(x^2+y^2)} \, dx\,dy .\]
The integrand is now radial, which is what makes polar coordinates pay. Writing $x = r\cos\theta$,
$y = r\sin\theta$, with Jacobian determinant $r$,
\[2I_\alpha = \int_0^{2\pi}\!\!\int_0^\infty r^2 e^{-\alpha r^2}\, r\,dr\,d\theta
  = 2\pi \int_0^\infty r^3 e^{-\alpha r^2}\,dr .\]
Substituting $u := r^2$, so that $du = 2r\,dr$, turns the remaining integral into one that can be
done by parts:
\[\int_0^\infty r^3 e^{-\alpha r^2}\,dr = \frac12 \int_0^\infty u e^{-\alpha u}\,du
  = \frac{1}{2\alpha^2}.\]
Therefore $2I_\alpha = 2\pi \cdot \frac{1}{2\alpha^2} = \frac{\pi}{\alpha^2}$, and
\[I_\alpha = \frac{\pi}{2\alpha^2}.\]
\begin{remark}
Two sanity checks come free. Scaling $(x,y) \mapsto (x,y)/\sqrt\alpha$ predicts $I_\alpha =
\alpha^{-2} I_1$, which the answer satisfies; and $I_\alpha$ should equal
$\tfrac12 \cdot \tfrac{1}{2\alpha}\cdot\tfrac{\pi}{\alpha}$ read off from the one-dimensional
Gaussian moments, which it does. The symmetry trick in the first line is the same one that
evaluates $\int_{\mathbb{R}} e^{-x^2}dx$ by squaring it: an integral that resists in Cartesian
coordinates becomes radial once it is paired with its own mirror image.
\end{remark}
\end{exercisesolution}


